% GENERATED FILE. Edit canonical lessons, then run npm run generate:books.
% canonical-source-hash: fnv1a64:248a2738462b62f2
% canonical-lessons: JA-C344-keikaku, JA-C344-youi-suru, JA-C344-youji, JA-C344-nittei, JA-C344-junchou

\chapter{Making Plans}
\label{ch:ja-making-plans}
\hlchaptermodality{344}

\begin{chapteropening}
\textbf{By the end of this chapter:} \emph{I can make a plan and say what happens if it changes.}

Complete the last lesson of chapter 344: I can make a plan and say what happens if it changes.
\end{chapteropening}

\section[\texorpdfstring{\ja{けいかく} (keikaku)}{けいかく (keikaku)}]{\ja{けいかく} (keikaku) — a scheme; a detailed plan}

\label{lesson:JA-C344-keikaku}



\begin{quote}
Before the new one: say the Japanese for to be found out, then the Japanese for to jump to conclusions.
\end{quote}

\subsection*{You'll want to know: \ja{けいかく}}
\textbf{\ja{けいかく}} — \emph{keikaku} — \textquotedblleft{}a scheme; a detailed plan\textquotedblright{}.

\begin{quote}
\textbf{\ja{もし} \ja{けいかくが} \ja{かわったら、れんらくして} \ja{ください。}} — \emph{moshi keikaku ga kawattara, renraku shite kudasai} — \textquotedblleft{}If the plan changes, please let me know.\textquotedblright{}
\end{quote}

\begin{cousinweb}
\textbf{\ja{けいかく}} (\emph{keikaku}) is Sino-Japanese (kango): two Chinese characters read with their on'yomi, readings borrowed from Chinese. They are \emph{kei}, \textquotedblleft{}plan\textquotedblright{} (Middle Chinese \emph{kejH}) and \emph{kaku}, \textquotedblleft{}brush-stroke\textquotedblright{} (\emph{hweak}) (KANJIDIC2; Middle Chinese in Baxter's transcription).
\end{cousinweb}

\begin{grammarlens}[title={What you've built}]
One more word for saying what happens if.
\end{grammarlens}

\subsection*{Your turn}
\begin{itemize}
\raggedright
  \item \emph{Say it:} \emph{keikaku}
  \item \emph{Say it:} \emph{keikaku}, once more
  \item \emph{Say it:} \emph{moshi keikaku ga kawattara, renraku shite kudasai} — If the plan changes, please let me know.
  \item \emph{Recall:} say the Japanese for to be found out, then the Japanese for to jump to conclusions, then say \emph{keikaku} again
\end{itemize}

\begin{tcolorbox}[breakable,colback=teal!4,colframe=teal!35!black,title={Before you move on}]
What does \ja{けいかく} mean? (\textquotedblleft{}A scheme; a detailed plan\textquotedblright{}.) Say it once more.
\end{tcolorbox}
\section[\texorpdfstring{\ja{よういする} (yōi suru)}{よういする (yōi suru)}]{\ja{よういする} (yōi suru) — to get things ready; to prepare}

\label{lesson:JA-C344-youi-suru}



\begin{quote}
Before the new one: say the Japanese for to jump to conclusions, then the Japanese for a detailed plan.
\end{quote}

\subsection*{You'll want to know: \ja{よういする}}
\textbf{\ja{よういする}} — \emph{yōi suru} — \textquotedblleft{}to get things ready; to prepare\textquotedblright{}.

\begin{quote}
\textbf{\ja{あめが} \ja{ふったら、かさを} \ja{よういして} \ja{ください。}} — \emph{ame ga futtara, kasa o yōi shite kudasai} — \textquotedblleft{}If it rains, please have an umbrella ready.\textquotedblright{}
\end{quote}

Polite \textbf{\ja{よういします}}, \emph{yōishimasu}; a \ja{する} verb. Past (polite): \textbf{\ja{ははが} \ja{おべんとうを} \ja{よういしました。}}, \emph{haha ga obentō o yōi shimashita}.

\begin{cousinweb}
\textbf{\ja{よういする}} (\emph{yōi suru}) is built from a Sino-Japanese noun and the native verb \emph{suru}, \textquotedblleft{}to do\textquotedblright{}. The noun, \emph{yōi}, is two Chinese characters read with their on'yomi, readings borrowed from Chinese: \emph{yō}, \textquotedblleft{}utilize\textquotedblright{} (Middle Chinese \emph{yowngH}) and \emph{i}, \textquotedblleft{}mind\textquotedblright{} (\emph{'iH}) (KANJIDIC2; JMdict; Middle Chinese in Baxter's transcription).
\end{cousinweb}

\begin{grammarlens}[title={What you've built}]
One more word for saying what happens if.
\end{grammarlens}

\subsection*{Your turn}
\begin{itemize}
\raggedright
  \item \emph{Say it:} \emph{yōi suru}
  \item \emph{Say it:} \emph{yōi suru}, once more
  \item \emph{Say it:} \emph{ame ga futtara, kasa o yōi shite kudasai} — If it rains, please have an umbrella ready.
  \item \emph{Recall:} say the Japanese for to jump to conclusions, then the Japanese for a detailed plan, then say \emph{yōi suru} again
\end{itemize}

\begin{tcolorbox}[breakable,colback=teal!4,colframe=teal!35!black,title={Before you move on}]
What does \ja{よういする} mean? (\textquotedblleft{}To get things ready; to prepare\textquotedblright{}.) Say it once more.
\end{tcolorbox}
\section[\texorpdfstring{\ja{ようじ} (yōji)}{ようじ (yōji)}]{\ja{ようじ} (yōji) — an errand; something to do}

\label{lesson:JA-C344-youji}



\begin{quote}
Before the new one: say the Japanese for a detailed plan, then the Japanese for to get things ready.
\end{quote}

\subsection*{You'll want to know: \ja{ようじ}}
\textbf{\ja{ようじ}} — \emph{yōji} — \textquotedblleft{}an errand; something to do\textquotedblright{}.

\begin{quote}
\textbf{\ja{ようじが} \ja{あったら、はやく} \ja{かえって} \ja{ください。}} — \emph{yōji ga attara, hayaku kaette kudasai} — \textquotedblleft{}If you have something to do, please go home early.\textquotedblright{}
\end{quote}

\begin{cousinweb}
\textbf{\ja{ようじ}} (\emph{yōji}) is Sino-Japanese (kango): two Chinese characters read with their on'yomi, readings borrowed from Chinese. They are \emph{yō}, \textquotedblleft{}utilize\textquotedblright{} (Middle Chinese \emph{yowngH}) and \emph{ji}, \textquotedblleft{}matter\textquotedblright{} (\emph{dzriH}) (KANJIDIC2; Middle Chinese in Baxter's transcription).
\end{cousinweb}

\begin{grammarlens}[title={What you've built}]
One more word for saying what happens if.
\end{grammarlens}

\subsection*{Your turn}
\begin{itemize}
\raggedright
  \item \emph{Say it:} \emph{yōji}
  \item \emph{Say it:} \emph{yōji}, once more
  \item \emph{Say it:} \emph{yōji ga attara, hayaku kaette kudasai} — If you have something to do, please go home early.
  \item \emph{Recall:} say the Japanese for a detailed plan, then the Japanese for to get things ready, then say \emph{yōji} again
\end{itemize}

\begin{tcolorbox}[breakable,colback=teal!4,colframe=teal!35!black,title={Before you move on}]
What does \ja{ようじ} mean? (\textquotedblleft{}An errand; something to do\textquotedblright{}.) Say it once more.
\end{tcolorbox}
\section[\texorpdfstring{\ja{にってい} (nittei)}{にってい (nittei)}]{\ja{にってい} (nittei) — a schedule (of days)}

\label{lesson:JA-C344-nittei}



\begin{quote}
Before the new one: say the Japanese for to get things ready, then the Japanese for an errand.
\end{quote}

\subsection*{You'll want to know: \ja{にってい}}
\textbf{\ja{にってい}} — \emph{nittei} — \textquotedblleft{}a schedule (of days)\textquotedblright{}.

\begin{quote}
\textbf{\ja{にっていが} \ja{きまったら、おしえて} \ja{ください。}} — \emph{nittei ga kimattara, oshiete kudasai} — \textquotedblleft{}When the dates are fixed, please tell me.\textquotedblright{}
\end{quote}

\begin{cousinweb}
\textbf{\ja{にってい}} (\emph{nittei}) is Sino-Japanese (kango): two Chinese characters read with their on'yomi, readings borrowed from Chinese. They are \emph{nichi}, \textquotedblleft{}day\textquotedblright{} (Middle Chinese \emph{nyit}) and \emph{tei}, \textquotedblleft{}extent\textquotedblright{} (\emph{drjeng}) (KANJIDIC2; Middle Chinese in Baxter's transcription). The two join as \emph{nit-tei}.
\end{cousinweb}

\begin{grammarlens}[title={What you've built}]
One more word for saying what happens if.
\end{grammarlens}

\subsection*{Your turn}
\begin{itemize}
\raggedright
  \item \emph{Say it:} \emph{nittei}
  \item \emph{Say it:} \emph{nittei}, once more
  \item \emph{Say it:} \emph{nittei ga kimattara, oshiete kudasai} — When the dates are fixed, please tell me.
  \item \emph{Recall:} say the Japanese for to get things ready, then the Japanese for an errand, then say \emph{nittei} again
\end{itemize}

\begin{tcolorbox}[breakable,colback=teal!4,colframe=teal!35!black,title={Before you move on}]
What does \ja{にってい} mean? (\textquotedblleft{}A schedule (of days)\textquotedblright{}.) Say it once more.
\end{tcolorbox}
\section[\texorpdfstring{\ja{じゅんちょう} (junchō)}{じゅんちょう (junchō)}]{\ja{じゅんちょう} (junchō) — going smoothly; favourable}

\label{lesson:JA-C344-junchou}



\begin{quote}
Before the new one: say the Japanese for an errand, then the Japanese for a schedule.
\end{quote}

\subsection*{You'll want to know: \ja{じゅんちょう}}
\textbf{\ja{じゅんちょう}} — \emph{junchō} — \textquotedblleft{}going smoothly; favourable\textquotedblright{}.

\begin{quote}
\textbf{\ja{はれたら、けいかくは} \ja{じゅんちょうでしょう。}} — \emph{haretara, keikaku wa junchō deshō} — \textquotedblleft{}If the weather clears, the plan will probably go smoothly.\textquotedblright{}
\end{quote}

A na-adjective: \ja{な} before a noun.

\begin{cousinweb}
\textbf{\ja{じゅんちょう}} (\emph{junchō}) is Sino-Japanese (kango): two Chinese characters read with their on'yomi, readings borrowed from Chinese. They are \emph{jun}, \textquotedblleft{}order\textquotedblright{} (Middle Chinese \emph{zywinH}) and \emph{chō}, \textquotedblleft{}tune\textquotedblright{} (\emph{dew}) (KANJIDIC2; Middle Chinese in Baxter's transcription).
\end{cousinweb}

\begin{grammarlens}[title={What you've built}]
One more word for saying what happens if.
\end{grammarlens}

\subsection*{Your turn}
\begin{itemize}
\raggedright
  \item \emph{Say it:} \emph{junchō}
  \item \emph{Say it:} \emph{junchō}, once more
  \item \emph{Say it:} \emph{haretara, keikaku wa junchō deshō} — If the weather clears, the plan will probably go smoothly.
  \item \emph{Recall:} say the Japanese for an errand, then the Japanese for a schedule, then say \emph{junchō} again
\end{itemize}

\begin{tcolorbox}[breakable,colback=teal!4,colframe=teal!35!black,title={Before you move on}]
What does \ja{じゅんちょう} mean? (\textquotedblleft{}Going smoothly; favourable\textquotedblright{}.) Say it once more.
\end{tcolorbox}
